\section{Building Blocks}
\label{sec:blocks}

The PLONK relation consists of the gate identity in
Eq.~\eqref{eq:plonk1} and the wiring permutation in
Eq.~\eqref{eq:plonk2}. By Corollary~\ref{corollary:schwartz},
we reduce the gate identity to a scalar claim by projecting onto
a random tensor vector
\(
    \ext{\bm{\xi}}
    =
    \tensorspan{\ext{\bm z}}
\)
and proving that the resulting inner product is zero. The projected
terms are inner products and weighted Hadamard products.
We use vector polynomial packing to reduce these two relations,
then use the reductions to prove permutation claims.

\subsection{Packed Inner-Product Proof}
\label{sec:ip}

Our first reduction packs a length-\(N\) inner-product claim over
\(\Field\) into a length-\(m\) claim over \(\ExtensionField\),
where \(N=m\kappa\). It plays the role of the degree-one zero
check in~\cite{iacr/DaoT24}.
Consider the relation:
\[
    c=\langle\bm a,\bm b\rangle,
    ~
    \bm a,\bm b\in\Field^N.
\]
Partition each vector into \(\kappa\) blocks of length \(m\), and
arrange these blocks as columns:
\begin{equation}
\label{eq:matAB}
    \bm A
    \coloneqq
    \concat{
        \bm a^{(1)},\ldots,\bm a^{(\kappa)}
    },
    \qquad
    \bm B
    \coloneqq
    \concat{
        \bm b^{(1)},\ldots,\bm b^{(\kappa)}
    }
    \in\Field^{m\times\kappa}.
\end{equation}
Encode the columns as vector polynomials:
\[
    \bm g_A(X)
    \coloneqq
    \sum_{j=1}^{\kappa}\bm a^{(j)}X^{j-1},
    \qquad
    \bm g_B(X)
    \coloneqq
    \sum_{j=1}^{\kappa}\bm b^{(j)}X^{j-1}
    \in\Field[X]^m,
\]
and pair evaluations at reciprocal points:
\begin{equation}
    f(X)
    \coloneqq
    \left\langle
        \bm g_A(X),\bm g_B(X^{-1})
    \right\rangle.
\label{eq:ip-laurent-polynomial}
\end{equation}
The constant coefficient of \(f\) is the sum of the matching-column
inner products, hence equals \(\langle\bm a,\bm b\rangle\).
All cross-column terms have nonzero exponents. Therefore,
\(c=\langle\bm a,\bm b\rangle\) if and only if there exist
\(h_+(X),h_-(X)\in\Field[X]\), each of degree at most
\(\kappa-2\), such that:
\begin{equation}
    f(X)
    =
    c
    +Xh_+(X)
    +X^{-1}h_-(X^{-1}).
\label{eq:ip-correction-identity}
\end{equation}

After the commitments to \(\bm A\) and \(\bm B\) are fixed, the
prover sends \(h_+\) and \(h_-\). The verifier then samples
\(
    \ext{x}\sample\ExtensionField^\times.
\)
The prover forms the packed vectors:
\[
    \ext{\bm a}
    \coloneqq
    \bm g_A(\ext{x}),
    \qquad
    \ext{\bm b}
    \coloneqq
    \bm g_B(\ext{x}^{-1})
    \in\ExtensionField^m.
\]
Evaluating Eq.~\eqref{eq:ip-correction-identity} at \(\ext{x}\)
reduces the original claim to:
\begin{equation}
    \left\langle
        \ext{\bm a},\ext{\bm b}
    \right\rangle
    =
    c
    +\ext{x}h_+(\ext{x})
    +\ext{x}^{-1}h_-(\ext{x}^{-1}).
\label{eq:packed-inner-product-check}
\end{equation}
The right-hand side is computable from the prover's two auxiliary
polynomials. Applying Fact~\ref{fact:ip} to the left-hand side
leaves two MLE claims on length-\(m\) packed vectors.
Algorithm~\ref{alg:ip} in Appendix~\ref{appx:pseudocode} gives
the protocol.

\begin{proposition}[Packed inner-product reduction]
\label{proposition:ip-complexity}
Algorithm~\ref{alg:ip} is complete and has reduction soundness
error at most:
\[
    \frac{2\kappa-2 }{|\ExtensionField|-1} + \frac{2\log_2 m}{|\ExtensionField|}.
\]
Excluding the terminal MLE openings, its prover cost is:
\[
    N(\kappa-1)\,\bb
    +
    2N\,\be
    +
    \left(
        \frac{4N}{\kappa}
        +
        O\left(
            \kappa+\log\frac{N}{\kappa}
        \right)
    \right)\ee.
\]
Its peak prover memory is \(O(N+\kappa^2)\) base-field
coordinates, which is \(O(N)\) whenever \(m\geq\kappa\).
\end{proposition}

We prove Proposition~\ref{proposition:ip-complexity} in
Appendix~\ref{appx:proof-ip}.

\subsection{Packed Hadamard-Product Proof}
\label{sec:hadamard}

We next consider a weighted Hadamard-product claim:
\begin{equation}
    \ext{c}
    =
    \left\langle
        \ext{\bm\xi},
        \bm a\circ\bm b
    \right\rangle,
    \qquad
    \ext{\bm\xi}
    =
    \ext{\bm r}\kronecker\ext{\bm l},
\label{eq:hadamard-claim}
\end{equation}
The vectors and weights have dimensions:
\[
    \bm a,\bm b\in\Field^N,
    \qquad
    \ext{\bm l}\in\ExtensionField^m,
    \qquad
    \ext{\bm r}\in\ExtensionField^\kappa,
    \qquad
    N=m\kappa.
\]
Using the block matrices \(\bm A,\bm B\in\Field^{m\times\kappa}\)
defined in Eq.~\eqref{eq:matAB}, Eq.~\eqref{eq:hadamard-claim}
becomes:
\begin{equation*}
    \ext{c}
    =
    \ext{\bm l}^{\top}
    (\bm A\circ\bm B)
    \ext{\bm r}.
\end{equation*}

The monomial encoding above extracts a sum over columns.
Here the columns have distinct weights, so we encode each column
claim at a separate interpolation node. The prover first sends:
\[
    \ext{\bm c}^{\top}
    \coloneqq
    \ext{\bm l}^{\top}
    (\bm A\circ\bm B)
    \in\ExtensionField^\kappa,
\]
whose entries are:
\[
    \ext{c}_j
    =
    \left\langle
        \ext{\bm l},
        \bm a^{(j)}\circ\bm b^{(j)}
    \right\rangle
    \qquad
    \text{for every }j\in[\kappa].
\]
The verifier checks:
\begin{equation*}
    \ext{c}
    \EqOrNot
    \left\langle
        \ext{\bm c},\ext{\bm r}
    \right\rangle.
\end{equation*}
It remains to verify the \(\kappa\) column claims together.

Let \(\{\psi_j(X)\}_{j=1}^{\kappa}\) and \(\Psi(X)\) be the
Lagrange basis and vanishing polynomial from
Definition~\ref{def:lagrange}. Define:
\[
\begin{aligned}
    \bm g_A(X)
    &\coloneqq
    \sum_{j=1}^{\kappa}
    \bm a^{(j)}\psi_j(X),
    \quad
    \bm g_B(X)
    \coloneqq
    \sum_{j=1}^{\kappa}
    \bm b^{(j)}\psi_j(X),
    \quad
    g_C(X)
    \coloneqq
    \sum_{j=1}^{\kappa}
    \ext{c}_j\psi_j(X).
\end{aligned}
\]
The coefficient fields are:
\[
    \bm g_A(X),\bm g_B(X)\in\Field[X]^m,
    \qquad
    g_C(X)\in\ExtensionField[X].
\]
At each interpolation node, they recover the corresponding column:
\[
    \bm g_A(\alpha_j)=\bm a^{(j)},
    \qquad
    \bm g_B(\alpha_j)=\bm b^{(j)},
    \qquad
    g_C(\alpha_j)=\ext{c}_j.
\]
Thus all column claims hold if and only if:
\[
    f(\alpha_j)=0
    \qquad
    \text{for every }j\in[\kappa],
\]
where:
\begin{equation*}
    f(X)
    \coloneqq
    \left\langle
        \ext{\bm l},
        \bm g_A(X)\circ\bm g_B(X)
    \right\rangle
    -
    g_C(X).
\end{equation*}
Since \(\deg f\leq2\kappa-2\), vanishing on \(\mathcal H\)
is equivalent to divisibility by \(\Psi\). Equivalently, there is
an \(h(X)\in\ExtensionField[X]\) of degree at most \(\kappa-2\)
such that:
\begin{equation}
    f(X)=\Psi(X)h(X).
\label{eq:hadamard-quotient}
\end{equation}

After the prover sends \(\ext{\bm c}\) and \(h\), the verifier
samples:
\[
    \ext{x}
    \sample
    \ExtensionField\setminus\mathcal H.
\]
The Lagrange evaluations give the packing weights:
\[
    \ext{\bm\rho}(\ext{x})
    \coloneqq
    \bigl(
        \psi_1(\ext{x}),\ldots,\psi_\kappa(\ext{x})
    \bigr)^\top.
\]
The prover computes the packed vectors:
\[
\begin{aligned}
    \ext{\bm a}
    &\coloneqq
    \bm A\ext{\bm\rho}(\ext{x})
    =
    \bm g_A(\ext{x}),
    \qquad
    \ext{\bm b}
    \coloneqq
    \bm B\ext{\bm\rho}(\ext{x})
    =
    \bm g_B(\ext{x}),
\end{aligned}
\]
and the verifier computes the interpolated claim:
\[
    \ext{c}_x
    \coloneqq
    g_C(\ext{x})
    =
    \left\langle
        \ext{\bm c},
        \ext{\bm\rho}(\ext{x})
    \right\rangle.
\]
Evaluating Eq.~\eqref{eq:hadamard-quotient} at \(\ext{x}\)
reduces the Hadamard-product claim to:
\begin{equation*}
    \left\langle
        \ext{\bm a},
        \ext{\bm b}\circ\ext{\bm l}
    \right\rangle
    =
    \ext{c}_x
    +
    \Psi(\ext{x})h(\ext{x}).
\end{equation*}
Fact~\ref{fact:ip} reduces this length-\(m\) inner product to two
MLE claims. Write its challenge point and terminal values as:
\[
    \ext{\bm y}\in\ExtensionField^{\log_2m},
    \qquad
    \ext{v}_1,\ext{v}_2\in\ExtensionField.
\]
The terminal claims on the original vectors are:
\begin{align}
\label{eq:hadamard-terminal}
    \ext{v}_1
    &\EqOrNot
    \left\langle
        \bm a,
        \ext{\bm\rho}(\ext{x})
        \kronecker
        \tensorspan{\ext{\bm y}}
    \right\rangle,
    \qquad
    \ext{v}_2
    \EqOrNot
    \left\langle
        \bm b,
        \ext{\bm\rho}(\ext{x})
        \kronecker
        \left(
            \ext{\bm l}
            \circ
            \tensorspan{\ext{\bm y}}
        \right)
    \right\rangle.
\end{align}

Algorithm~\ref{alg:hadamard} in Appendix~\ref{appx:pseudocode}
summarizes the protocol.

\begin{proposition}[Packed Hadamard-product proof]
\label{proposition:hadamard-complexity}
Algorithm~\ref{alg:hadamard} is complete. With the terminal MLE
claims verified, its soundness error is at most:
\[
\begin{aligned}
    \epsilon_{\mathsf{HP}}
    &\leq
    \frac{2\kappa-2}
         {|\ExtensionField|-\kappa}
    +
    \frac{2\log_2N}
         {|\ExtensionField|}.
\end{aligned}
\]
For \(\mathcal H\subseteq\Field\), with precomputed coefficients
for the Lagrange basis and \(\Psi\), the dominant prover cost
under schoolbook arithmetic is:
\[
    3\kappa N\,\bb
    +
    7N\,\be
    +
    \frac{6N}{\kappa}\,\ee.
\]
More precisely, the prover cost is:
\[
\begin{aligned}
    &(3\kappa+1)N\,\bb+
    \left(
        7N-\frac{N}{\kappa}
        +O(\kappa^2)
    \right)\be+
    \left(
        \frac{6N}{\kappa}
        +
        O\left(
            \kappa+\log\frac{N}{\kappa}
        \right)
    \right)\ee.
\end{aligned}
\]
The peak prover memory is \(O(N+\kappa^2)\) base-field
coordinates, which is \(O(N)\) whenever \(m\geq\kappa\).
\end{proposition}

We prove Proposition~\ref{proposition:hadamard-complexity} in
Appendix~\ref{appx:proof-hadamard}.

\smallskip\noindent\textbf{Binary-field variant.}
For \(\Field=\mathbb F_2\) and \(\kappa=2^s\), Boolean
tensor interpolation replaces the univariate quotient by \(s\)
quotients for the generators \(X_j^2-X_j\).
Given \(\ext{\bm l}\), all quotients can be constructed in a
mixed tensor basis using \(O(m\kappa^{\log_2 3})\) bit operations
and the same order of extension-field XORs, without general
extension-field multiplication in this stage.
Appendix~\ref{appx:binary-hadamard} gives the construction,
separate bit and word costs, and the subsequent identity test.
Proposition~\ref{proposition:binary-hadamard-complexity}
summarizes the total complexity and soundness, including the
packed-opening reductions.

\smallskip\noindent\textbf{Small fields via an intermediate field.}
For \(q=|\Field|<\kappa\), we interpolate over an intermediate
field. With \(\kappa\) a power of two, choose \(t\)
as the smallest power of two satisfying \(q^t\geq2\kappa-1\).
Then \(t\mid\kappa\) and \(t=O(\log_q\kappa)\), giving the tower:
\[
    \Field\subseteq\mathbb K=\mathbb F_{q^t}
    \subseteq\ExtensionField,
    \qquad [\ExtensionField:\mathbb K]=\kappa/t.
\]
Choose \(\mathcal H\subseteq\mathbb K\) of size \(\kappa\),
define \(\bm\psi\) and \(\Psi\) as above, and take
\(\mathcal S=\{\zeta_1,\ldots,\zeta_{2\kappa-1}\}\subseteq\mathbb K\)
containing \(\mathcal H\).
Set \(P_{j\ell}=\psi_j(\zeta_\ell)\).
The witness matrices remain over \(\Field\).
In a fixed \(\mathbb K\)-basis \(\ext{\bm\eta}\) of
\(\ExtensionField\), write \(\ext{\bm l}=\bm L\ext{\bm\eta}\),
where \(\bm L\in\mathbb K^{m\times(\kappa/t)}\).
Compute the aggregate samples:
\begin{equation}
    \bm U
    =\bm L^\top\bigl((\bm A\bm P)\circ(\bm B\bm P)\bigr)
    \in\mathbb K^{(\kappa/t)\times(2\kappa-1)}.
\label{eq:hadamard-small-field-aggregation}
\end{equation}
Interpolating \(\bm U^\top\ext{\bm\eta}\), subtracting \(g_C\),
and dividing by \(\Psi\) recovers \(h\) of degree at most
\(\kappa-2\). Thus interpolation is performed only after aggregation.

With field representations, interpolation constants, and basis
maps precomputed, schoolbook arithmetic gives the following
costs, excluding commitments and openings:
\begin{equation}
    T_{\Prover}=O\bigl((N\kappa+\kappa^3)t\bigr),
    \qquad S_{\Prover}=O(N+\kappa^2).
\label{eq:hadamard-small-field-cost}
\end{equation}
Time and space are measured in base-field operations and coordinates.
The space bound follows by generating and aggregating one sample
column at a time; Appendix~\ref{appx:hadamard-comparison} gives
the counts. For \(m\geq\kappa\), the bounds simplify to
\(O(N\kappa\log_q\kappa)\) time and \(O(N)\) space.

\smallskip\noindent\textbf{Tabulated implementation.}
Two types of tables remove the factor \(t\).
First, tabulate multiplication in \(\mathbb K\), including
\(\mathbb K\times\Field\).
Charging table accesses per base-field coordinate, each product
costs \(O(t)\). Aggregation then costs \(O(N\kappa)\), and
interpolation and division cost \(O(\kappa^3)\).

Second, use blockwise linear-map tables, as in Flock's binary
LDE implementation~\cite[Section~4.2]{iacr/BunzRW26}, to tabulate
sample formation in blocks of
\(\tau=\lfloor\log_q\kappa\rfloor\), so that
\(q^\tau\leq\kappa\) and \(t/\tau=O(1)\).
Partition the witness columns and corresponding rows of \(\bm P\)
into blocks of size \(\tau\), padding the last block with zeros.
For each block \(\bm P^{(j)}\), store:
\begin{equation}
    \mathcal T_j[\bm v]=\bm v^\top\bm P^{(j)},
    \qquad \bm v\in\Field^\tau.
\label{eq:hadamard-block-table}
\end{equation}
Each output row sums \(\lceil\kappa/\tau\rceil\) entries of
\(O(\kappa t)\) base-field coordinates. Hence forming both
\(\bm A\bm P\) and \(\bm B\bm P\), including all reads and
additions, costs \(O(N\kappa t/\tau)=O(N\kappa)\).
The block tables occupy
\(O(\lceil\kappa/\tau\rceil q^\tau\kappa t)=O(\kappa^3)\)
coordinates and can be built in \(O(\kappa^3)\) operations:
precompute scalar multiples of the rows of \(\bm P^{(j)}\)
and extend partial input assignments one coordinate at a time.

The online cost and total storage are therefore:
\begin{equation}
\begin{aligned}
    T_{\Prover}&=O(N\kappa+\kappa^3),\\
    S_{\Prover}&=O(N+\kappa^3+tq^{2t}).
\end{aligned}
\label{eq:hadamard-block-tabulated-cost}
\end{equation}
The \(q^{2t}\)-entry multiplication table and the block tables
require \(O(t^2q^{2t}+\kappa^3)\) preprocessing operations.
Storage includes all read-only tables; processing witness rows
one at a time uses \(O(N+\kappa^2)\) working coordinates.
For \(m\geq\kappa\), the online cost is \(O(N\kappa)\).

The quotient degree, challenge order, and terminal queries in
Eq.~\eqref{eq:hadamard-terminal} are unchanged. Hence the
soundness bound of Proposition~\ref{proposition:hadamard-complexity}
and the packed-opening interface of Section~\ref{sec:pcs} still
apply; the costs are those in
Eq.~\eqref{eq:hadamard-block-tabulated-cost}.

\subsection{Packed Permutation Proof}
\label{sec:permutation}

We use packed Hadamard proofs to establish the permutation relation:
\[
    \bm a\stackrel{\mathsf{perm}}{=}\bm b,
    \qquad
    \bm a,\bm b\in\Field^N.
\]
After the commitments to \(\bm a\) and \(\bm b\) are fixed, the
verifier samples
\(
    \ext{z}\sample\ExtensionField
\)
and checks equality of the two product fingerprints:
\begin{equation}
   \Theta \coloneqq  \mathsf{GrandProd}
    \bigl(\ext{z}\bm 1+\bm a\bigr)
    \EqOrNot
    \mathsf{GrandProd}
    \bigl(\ext{z}\bm 1+\bm b\bigr).
\label{eq:packed-permutation-fingerprint}
\end{equation}
Permutations have identical fingerprints at every \(\ext z\).
Otherwise, consider the difference polynomial:
\[
    \prod_{i=1}^{N}(Z+a_i)
    -
    \prod_{i=1}^{N}(Z+b_i).
\]
It is nonzero and has degree at most \(N-1\). Therefore,
Eq.~\eqref{eq:packed-permutation-fingerprint} holds for a
uniformly random \(\ext z\) with probability at most
\(
    \min\left\{
        1,\frac{N-1}{|\ExtensionField|}
    \right\}.
\)

It therefore suffices to prove each grand product in
Eq.~\eqref{eq:packed-permutation-fingerprint}. We describe the
proof for one committed vector and apply it twice when both
vectors are committed. A public vector's grand product can be
computed directly.

Let \(\bm d\in\Field^N\) denote the vector whose grand product is
to be proved. Assume:
\[
    N=m\kappa,
    \qquad
    m,\kappa\in2^{\mathbb N},
\]
and arrange its \(\kappa\) length-\(m\) blocks as:
\[
    \bm D
    \coloneqq
    \concat{
        \bm d^{(1)},\ldots,\bm d^{(\kappa)}
    }
    \in\Field^{m\times\kappa}.
\]
The two vectors may be padded identically to meet these dimensions.

We first multiply across columns over \(\Field[Z]\), leaving
\(m\) row products to evaluate over \(\ExtensionField\).
Define the column-product tree:
\begin{equation}
\begin{aligned}
    \bm C_0(Z)
    &\coloneqq
    Z\bm 1_{m\times\kappa}+\bm D,\\
    \bm C_t(Z)
    &=
    \concat{
        \bm C_{t,L}(Z),
        \bm C_{t,R}(Z)
    },\\
    \bm C_{t+1}(Z)
    &\coloneqq
    \bm C_{t,L}(Z)
    \circ
    \bm C_{t,R}(Z),
    \qquad
    t\in\{0,\ldots,\log_2\kappa-1\}.
\end{aligned}
\label{eq:packed-permutation-product-tree}
\end{equation}
At level \(t\):
\[
    \bm C_t(Z)
    \in
    \Field[Z]^{m\times(\kappa/2^t)}
    \qquad\text{and}\qquad
    \deg_Z\bm C_t\leq2^t.
\]
The final layer is a single polynomial column:
\[
    \bm C_{\log_2 \kappa}(Z)
    \in\Field[Z]^m,
\]
satisfying:
\begin{equation*}
    \mathsf{GrandProd}
    \bigl(\ext z\bm 1+\bm d\bigr)
    =
    \mathsf{GrandProd}
    \bigl(\bm C_{\log_2 \kappa}(\ext z)\bigr).
\end{equation*}

To prove the grand product, we traverse two trees in reverse:
an extension-field tree over the \(m\) evaluated row products,
followed by the polynomial tree across the \(\kappa\) columns.
For the first tree, let
\(
    \ext{\bm c}
    \coloneqq
    \bm C_{\log_2\kappa}(\ext z)
    \in\ExtensionField^m
\)
through a sequence of recursive Hadamard-product proofs. Set
\(
    \ext{\bm c}_0
    \coloneqq
    \ext{\bm c}.
\)
For each \(i=0,\ldots,\log_2m-1\), split
\(
    \ext{\bm c}_i
    =
    \ext{\bm c}_{i,L}
    \mathbin{\|}
    \ext{\bm c}_{i,R}
\)
into equal halves and define:
\begin{equation*}
    \ext{\bm c}_{i+1}
    \coloneqq
    \ext{\bm c}_{i,L}
    \circ
    \ext{\bm c}_{i,R}.
\end{equation*}
Thus,
\(
    \ext{\bm c}_i
    \in
    \ExtensionField^{m/2^i},
\)
and the final layer consists of the single scalar
\(
    \ext c_{\log_2m}
    =
    \Theta
    =
    \mathsf{GrandProd}(\ext{\bm c})
    =
    \mathsf{GrandProd}
    \bigl(\ext z\bm 1+\bm d\bigr).
\)
Each transition is a Hadamard-product relation:
\[
    \ext{\bm c}_{i+1}
    =
    \ext{\bm c}_{i,L}
    \circ
    \ext{\bm c}_{i,R}.
\]
Starting from the claimed product \(\Theta\), the protocol
reduces these relations from \(i=\log_2m-1\) down to \(i=0\).

The remaining claim concerns the evaluated terminal column:
\begin{equation*}
    \ext c_{\log_2 \kappa}
    \EqOrNot
    \left\langle
        \ext{\bm\xi}_{\log_2 \kappa},
        \bm C_{\log_2 \kappa}(\ext z)
    \right\rangle.
\end{equation*}

We then unfold the column-product tree from
\(t=\log_2\kappa-1\) down to \(t=0\). Suppose the current
claim is:
\begin{equation*}
    \ext c_{t+1}
    \EqOrNot
    \left\langle
        \ext{\bm\xi}_{t+1},
        \operatorname{vec}
        \bigl(\bm C_{t+1}(\ext z)\bigr)
    \right\rangle.
\end{equation*}
By Eq.~\eqref{eq:packed-permutation-product-tree}, this is the
weighted Hadamard-product relation:
\begin{equation}
\begin{aligned}
    \ext c_{t+1}
    \EqOrNot
    \Bigl\langle
        \ext{\bm\xi}_{t+1},\,
        &\operatorname{vec}
        \bigl(\bm C_{t,L}(\ext z)\bigr)
        \circ\,
        \operatorname{vec}
        \bigl(\bm C_{t,R}(\ext z)\bigr)
    \Bigr\rangle.
\end{aligned}
\label{eq:packed-permutation-hadamard-claim}
\end{equation}

Represent each entry of \(\bm C_t(Z)\) by its base-field
coefficients and absorb evaluation at \(\ext z\) into the
tensor weight in
Eq.~\eqref{eq:packed-permutation-hadamard-claim}. The prover and
verifier apply Algorithm~\ref{alg:hadamard} and randomly combine
the two resulting MLE claims into one claim for the preceding layer:
\begin{equation*}
    \ext c_t
    \EqOrNot
    \left\langle
        \ext{\bm\xi}_t,
        \operatorname{vec}
        \bigl(\bm C_t(\ext z)\bigr)
    \right\rangle.
\end{equation*}
After \(\log_2\kappa\) reductions, this gives:
\[
    \ext c_0
    \EqOrNot
    \left\langle
        \ext{\bm\xi}_0,
        \operatorname{vec}
        \bigl(\bm C_0(\ext z)\bigr)
    \right\rangle.
\]
Since:
\[
    \bm C_0(\ext z)
    =
    \ext z\bm 1_{m\times\kappa}+\bm D,
\]
the final claim is an opening of the original vector:
\begin{equation*}
    \left\langle
        \ext{\bm\xi}_0,\bm d
    \right\rangle
    \EqOrNot
    \ext c_0
    -
    \ext z
    \left\langle
        \ext{\bm\xi}_0,\bm 1_N
    \right\rangle.
\end{equation*}
Applying this reduction to both vectors and checking
Eq.~\eqref{eq:packed-permutation-fingerprint} proves the permutation
claim. Algorithm~\ref{alg:permutation} in
Appendix~\ref{appx:pseudocode} gives the protocol.

\begin{proposition}[Packed permutation proof]
\label{proposition:permutation-complexity}
The packed permutation proof is complete. Let
\(\epsilon_{\mathsf{GP}}\) denote the soundness error of proving
one grand product, excluding the terminal PCS openings. Then:
\[
\begin{aligned}
    \epsilon_{\mathsf{GP}}
    \leq\;&
    \frac{
        \log_2m(\log_2m-1)
    }{
        |\ExtensionField|
    }
    +
    \frac{
        2\log_2m\log_2\kappa
    }{
        |\ExtensionField|
    }
    +
    \sum_{t=0}^{\log_2\kappa-1}
    \frac{
        2^{\log_2\kappa-t+1}-2
    }{
        |\ExtensionField|
        -
        2^{\log_2\kappa-t}
    }.
\end{aligned}
\]
Consequently, when both \(\bm a\) and \(\bm b\) are committed,
the permutation proof has soundness error at most:
\[
    \min\left\{
        1,
        \frac{N-1}{|\ExtensionField|}
    \right\}
    +
    2\epsilon_{\mathsf{GP}}
    +
    \epsilon_{\mathsf{PCS}},
\]
where \(\epsilon_{\mathsf{PCS}}\) denotes the soundness error of
the terminal MLE openings.

Under schoolbook arithmetic, the prover cost for one
grand-product proof is:
\[
    O(N\kappa)\,\bb
    +
    O(N)\,\be
    +
    \left(
        \frac{10N}{\kappa}
        +
        O\left(
            \kappa
            +
            (\log_2m)^2
            +
            \log_2m\log_2\kappa
        \right)
    \right)\ee.
\]
The peak prover memory is \(O(N+\kappa^2)\) base-field
coordinates, which is \(O(N)\) whenever \(m\geq\kappa\).
Proving both committed grand products doubles the corresponding
prover costs.
\end{proposition}

We prove Proposition~\ref{proposition:permutation-complexity} in
Appendix~\ref{appx:proof-permutation}.

\subsection{Packed Polynomial Commitment Scheme}
\label{sec:pcs}

The preceding reductions query a committed matrix using packing
weights chosen after the commitment.
Let \(\bm A\in\Field^{m\times\kappa}\), with \(N=m\kappa\),
and fix an \(\Field\)-basis \(\ext{\bm\beta}\) of
\(\ExtensionField\).
We commit to a fixed-basis packing and authenticate weighted
evaluations:
\[
    \ext{\bm a}=\bm A\ext{\bm\beta},
    \qquad
    \ext u=\langle\bm A\ext{\bm\rho},\ext{\bm\chi}\rangle,
\]
where \(\ext{\bm\rho}\in\ExtensionField^\kappa\) and
\(\ext{\bm\chi}\in\ExtensionField^m\) is given by its two-entry
tensor factors.

We use ring-switching~\cite{eurocrypt/DiamondP26}, expressed in
our vector notation, to reduce each such claim to one MLE opening
of \(\ext{\bm a}\).
Coordinate decomposition and a random projection give an
extension-field inner-product claim, while the verifier evaluates
the induced public query directly from its tensor factors.
Under schoolbook arithmetic and \(m\geq\kappa\), the prover uses
\(O(\kappa N)\) base-field operations and \(O(N)\) base-field
coordinates of memory for the opening reduction, excluding the
underlying PCS.
Appendix~\ref{appx:pcs} gives the construction and its analysis,
including an instantiation from a base-field linear code via
Blaze's code-switching construction~\cite{eurocrypt/BrehmCFRRZ25}.
